%% ============================================================
%%  Companion Handout -- Bayesian Estimation from First Principles
%%  Standalone teaching document (NOT part of the book's Part I).
%%  Written by the student as a from-zero introduction to the
%%  material of Lecture 2, SS2.8-2.9; self-contained, no statistics
%%  background assumed. Compile standalone: pdflatex handout-bayes-estimation.
%% ============================================================
\documentclass[main.tex]{subfiles}

\begin{document}

\beginunit{Handout~\textendash~Bayesian Estimation}{COMPANION HANDOUT}%
  {Written for the course notes \textperiodcentered\ Summer 2026}%
  {Bayesian Estimation from First Principles}

\epigraph{Probability theory is nothing but common sense reduced to calculation.}%
{Pierre-Simon Laplace, \emph{Essai philosophique sur les probabilit\'es} (1814)}

\begin{lecturecontents}{Contents of this handout}
  \item \hyperlink{B.1}{B.1\quad The Problem: Guessing a Number You Cannot See}
  \item \hyperlink{B.2}{B.2\quad The Probability Toolkit, from Zero}
  \item \hyperlink{B.3}{B.3\quad Estimation as a Decision, and Why It Gets Stuck}
  \item \hyperlink{B.4}{B.4\quad The Prior, the Posterior, and Bayes' Theorem}
  \item \hyperlink{B.5}{B.5\quad The Best Guess Under Squared Error}
  \item \hyperlink{B.6}{B.6\quad A Toolkit Distribution: the Beta Family}
  \item \hyperlink{B.7}{B.7\quad The Cure Rate, Start to Finish}
  \item \hyperlink{B.8}{B.8\quad A Second Instrument: the Normal--Normal Pair}
  \item \hyperlink{B.9}{B.9\quad What to Make of the Prior}
  \item \hyperlink{B.10}{B.10\quad The Vocabulary, and Where to Go Next}
\end{lecturecontents}

\bigskip
\noindent\textbf{About this handout.} This is a student-written companion to the course
notes---not an instructor document. It rebuilds, from nothing, the mathematics behind one
sentence from Lecture~2: \emph{under squared-error loss, the best estimator is the posterior
mean}. Every term is defined before it is used, every derivation shows every algebraic step,
and every step says why it is being taken. No statistics background is assumed; the only
prerequisites are single-variable calculus (derivatives, integrals, integration by parts) and
patience. Cross-references like \xrefL{2.9.3} point at numbered objects in the course notes
(lecture~2, section~9, object~3); references like \xref{B.5.1} point within this handout.
A web edition lives on the course site, identical in content, with the figures replaced by
live laboratories---the normal--normal update of \S\,\xref{B.8} as the centerpiece---so the
machinery can be steered by hand while reading.

\clearpage
% ============================================================
\sub{B.1}{The Problem: Guessing a Number You Cannot See}

\begin{ObjL}{Concept}{B.1.1}{The estimation problem}
A new treatment cures some fraction of the patients who receive it. Call that fraction
$\theta$: if $\theta=0.7$, then in the long run $70\%$ of treated patients are cured. Nobody
knows $\theta$. It is not random---it is a fixed fact about the treatment---but it is
\emph{hidden}: the only way to learn about it is to treat patients and watch what happens.
Suppose we treat $n=10$ patients and $x=7$ of them are cured. What number should we now
report as our best guess of $\theta$? That question---turn data into a single defensible
guess at a hidden number---is called \emph{estimation}, and this handout builds one complete,
principled answer to it.
\end{ObjL}

\begin{ObjL}{Concept}{B.1.2}{Two sources of information}
The obvious answer is the observed cure rate, $7/10=0.7$. It is a good answer, and by the end
of the handout you will know exactly in what sense. But notice that the data are not the only
thing you know. Perhaps treatments of this kind have historically cured between $30\%$ and
$60\%$ of patients; perhaps the biology says a cure rate near $1$ is implausible. Ten
patients is not many, and it seems wasteful to ignore everything except those ten outcomes.
The \emph{Bayesian} approach--- the subject of this handout---is defined by one move: it
expresses \emph{both} sources of information, the data and the background knowledge, in the
same language, the language of probability. Once both speak that language, a single theorem
combines them optimally.
\end{ObjL}

\begin{ObjL}{Remark}{B.1.3}{The route}
The handout runs in a straight line, and it helps to see the whole road first.
\S\,\xref{B.2} builds the probability vocabulary (random variables, distributions,
expectation, conditional distributions)---the language everything else is written in.
\S\,\xref{B.3} makes ``estimation'' precise and shows why, without background knowledge, the
problem gets \emph{stuck}: candidate guessing rules cannot be ranked. \S\,\xref{B.4}
introduces the Bayesian move---treat the hidden number as random, with a distribution
expressing what you believe about it---and derives \emph{Bayes' theorem}, the rule for
updating that belief with data. \S\,\xref{B.5} proves the punchline: once beliefs are
updated, the best single-number guess under the natural error measure is the \emph{average
of the updated belief}. \S\S\,\xref{B.6}--\xref{B.8} then work two complete examples with
real numbers, including the ten-patient trial above, and \S\,\xref{B.9} asks the honest
question---where did the prior belief come from?---before \S\,\xref{B.10} hands you the
standard vocabulary for further reading.
\end{ObjL}

% ============================================================
\sub{B.2}{The Probability Toolkit, from Zero}

Probability is the mathematics of quantities whose values you cannot predict. This section
defines the six ideas the rest of the handout uses---random variable, distribution,
expectation, variance, conditional distribution, conditional expectation---each with a small
computation so the definition is not just words. Readers who have met these before can skim,
but should slow down at \xref{B.2.7} and \xref{B.2.9}, which carry the whole later argument.

\begin{Obj}{Definition}{B.2.1}{Random variable; realization}
A \emph{random variable} is a numerical quantity whose value is determined by the outcome of
something uncertain---a measurement not yet made. The number of cures among ten patients,
before the trial runs, is a random variable; so is tomorrow's temperature. Random variables
are written with capital letters ($X$, $Y$); once the uncertainty resolves and a value is
observed, that observed number---written in lowercase ($x$, $y$)---is called the
\emph{realization}. ``$X=7$'' is an event that had some probability in advance; ``$x=7$'' is
what actually happened.
\end{Obj}

\begin{Obj}{Definition}{B.2.2}{Distribution; pmf; density}
The \emph{distribution} of a random variable is the complete advance description of how
likely it is to fall anywhere: the assignment of a probability to every statement of the
form ``$X$ lands in the set $A$.'' For a \emph{discrete} random variable (one taking
isolated values, like a count) the distribution is carried by the \emph{probability mass
function} (pmf)
\[ p(x) \;=\; P(X = x) , \]
the probability of each individual value; probabilities of sets are found by adding:
$P(X\in A)=\sum_{x\in A}p(x)$. For a \emph{continuous} random variable (one taking values in
a continuum, like a fraction in $[0,1]$) individual points all have probability zero, and
the distribution is instead carried by a \emph{density} $f(x)\ge0$, with probabilities found
by integrating:
\[ P(X\in A) \;=\; \int_A f(x)\,dx .\note{A density is not a probability but a probability
\emph{rate}: probability per unit length, so that $f(x)\,dx\approx P(x\le X\le x+dx)$ for
small $dx$. Consequently a density can exceed $1$ (the uniform density on $[0,\frac12]$ is
$2$ there) and $f(x)$ is emphatically not $P(X=x)$, which is $0$. Sums against a pmf and
integrals against a density are the same idea with the same rules; the handout states things
once and trusts the reader to swap $\sum\leftrightarrow\int$.} \]
In both cases the total is $1$: $\sum_x p(x)=1$, $\int f(x)\,dx=1$---\emph{something} must
happen.
\end{Obj}

\begin{Obj}{Definition}{B.2.3}{Expectation}
The \emph{expectation} (or \emph{mean}) of a random variable $X$ is its
probability-weighted average value,
\[ \E\{X\} \;=\; \sum_x x\, p(x) \qquad\text{(discrete)}, \qquad\qquad
   \E\{X\} \;=\; \int x\, f(x)\, dx \qquad\text{(continuous)} : \]
each possible value, weighted by how likely it is. Geometrically it is the \emph{balancing
point} of the distribution: if the pmf or density were a physical distribution of weight
along a rod, $\E\{X\}$ is where the rod balances. It is the single number that best
summarizes ``where the distribution sits''---a claim \S\,\xref{B.5} will turn from slogan
into theorem.
\end{Obj}

\begin{ObjL}{Example}{B.2.4}{Two expectations, fully computed}
(a) A random variable $Z$ takes the value $1$ with probability $\theta$ and $0$ with
probability $1-\theta$ (one patient: cured or not; such a $Z$ is called a
\emph{Bernoulli($\theta$)} variable). Its expectation, straight from the definition:
\[ \E\{Z\} \;=\; 1\cdot P(Z=1) + 0\cdot P(Z=0) \;=\; 1\cdot\theta + 0\cdot(1-\theta)
   \;=\; \theta . \]
The mean of a $0$-or-$1$ variable \emph{is} its probability of being $1$---the bridge
between ``probability'' and ``long-run fraction.''
(b) A random variable $W$ takes values $1, 2, 6$ with probabilities
$\tfrac12, \tfrac13, \tfrac16$. Then
\[ \E\{W\} \;=\; 1\cdot\tfrac12 + 2\cdot\tfrac13 + 6\cdot\tfrac16
   \;=\; \tfrac12 + \tfrac23 + 1 \;=\; \tfrac{3+4+6}{6} \;=\; \tfrac{13}{6}
   \;\approx\; 2.17 . \]
Note the balancing-point reading: $2.17$ is not a possible value of $W$; it is where the
weights balance. Expectation also respects sums and constant multiples---%
$\E\{aX+bY\}=a\,\E\{X\}+b\,\E\{Y\}$, always, with no assumptions about how $X$ and $Y$ are
related---a property called \emph{linearity} that we will use constantly.\note{Why linearity
is free: in $\E\{aX+bY\}=\sum (ax+by)\,p(x,y)$ the sum splits into
$a\sum x\,p(x,y)+b\sum y\,p(x,y)$ by ordinary algebra, and the two pieces are $a\,\E\{X\}$
and $b\,\E\{Y\}$. Nothing about the relationship between $X$ and $Y$ was used---only that
both means exist.}
\end{ObjL}

\begin{Obj}{Definition}{B.2.5}{Variance; standard deviation}
The \emph{variance} of $X$ measures its spread: the expected squared distance from its own
mean,
\[ \Var(X) \;=\; \E\big\{ (X - \E\{X\})^2 \big\} . \]
Squaring makes deviations positive (so left- and right-misses cannot cancel) and charges
large deviations disproportionately. Its square root, the \emph{standard deviation}, is the
spread in the original units. A useful equivalent form, obtained by expanding the square and
using linearity:
\[ \Var(X) \;=\; \E\big\{X^2 - 2X\,\E\{X\} + \E\{X\}^2\big\}
   \;=\; \E\{X^2\} - 2\,\E\{X\}\,\E\{X\} + \E\{X\}^2 , \]
that is,
\[ \Var(X) \;=\; \E\{X^2\} - \E\{X\}^2 . \]
For the Bernoulli variable of \xref{B.2.4}(a): $Z^2=Z$ (since $0^2=0$, $1^2=1$), so
$\E\{Z^2\}=\theta$ and
\[ \Var(Z) \;=\; \theta - \theta^2 \;=\; \theta(1-\theta) , \]
largest at $\theta=\tfrac12$ (a fair coin is the most unpredictable coin) and zero at
$\theta=0$ or $1$ (no uncertainty at all).
\end{Obj}

\begin{Obj}{Definition}{B.2.6}{Joint, marginal, and conditional distributions}
For a \emph{pair} of random variables $(X,Y)$, the \emph{joint distribution} assigns a
probability to every combination: in the discrete case, $p(x,y)=P(X=x,\,Y=y)$. From it, the
distribution of each variable alone---the \emph{marginal}---is recovered by summing out the
other:
\[ p_X(x) \;=\; \sum_y p(x,y) . \]
And the \emph{conditional distribution} of $X$ given that $Y$ has been observed to equal $y$
is the joint, restricted to that observation and rescaled to total $1$:
\[ p(x \mid y) \;=\; \frac{p(x,y)}{p_Y(y)} \qquad (p_Y(y)>0) . \]
Read the formula as a two-step act: \emph{take the slice} of the joint where $Y=y$ (the
numerator), then \emph{renormalize} it into a genuine probability distribution by dividing
by the slice's total mass (the denominator). For continuous variables the same formulas hold
with densities and integrals.
\end{Obj}

\begin{ObjL}{Example}{B.2.7}{A worked table---and, secretly, the whole handout}
A disease affects $10\%$ of a population. A screening test comes back positive for $90\%$ of
the sick and (falsely) for $20\%$ of the healthy. Let $D=1$ mean sick, $T=1$ mean the test is
positive. The joint distribution, built by multiplying each group's size by its test
behavior:
\[ \begin{array}{c|cc|c}
     & T=1 & T=0 & \text{row total}\\ \hline
   D=1 & 0.10\times0.90 = 0.09 & 0.10\times0.10 = 0.01 & 0.10\\
   D=0 & 0.90\times0.20 = 0.18 & 0.90\times0.80 = 0.72 & 0.90\\ \hline
   \text{column total} & 0.27 & 0.73 & 1.00
   \end{array} \]
The four cells total $1$, the row totals are the marginal of $D$, the column totals the
marginal of $T$. Now suppose your test comes back positive. What is the probability you are
sick? Apply \xref{B.2.6}: slice the $T=1$ column, then renormalize by its total $0.27$:
\[ P(D=1 \mid T=1) \;=\; \frac{p(1,1)}{p_T(1)} \;=\; \frac{0.09}{0.27} \;=\; \frac13 . \]
Pause on what just happened. \emph{Before} the test, your probability of being sick was
$0.10$. You received data (a positive test), and your probability \emph{updated} to
$0.33$---higher, but nowhere near certainty, because false positives are common and the
disease is rare. This little table is the entire Bayesian method in miniature: a belief, a
data-generating mechanism, an observation, an updated belief. \S\,\xref{B.4} will redo this
computation with a continuous unknown in place of $D$ and give its parts their proper names.
\end{ObjL}

\begin{Obj}{Definition}{B.2.8}{Conditional expectation}
The \emph{conditional expectation} of $X$ given $Y=y$ is the mean of the conditional
distribution:
\[ \E\{X \mid Y=y\} \;=\; \sum_x x\, p(x\mid y) \]
(integral against the conditional density in the continuous case). It is the balancing point
of the \emph{slice}: your updated best summary of $X$ after seeing $Y=y$. Note that it is a
number for each $y$---so, before $Y$ is observed, the recipe $y\mapsto\E\{X\mid Y=y\}$ is
itself a random quantity, written $\E\{X\mid Y\}$.
\end{Obj}

\begin{Obj}{Theorem}{B.2.9}{Iterated expectations}
For jointly distributed $(X,Y)$ with $\E\{|X|\}<\infty$,
\[ \E\{X\} \;=\; \E\big\{\, \E\{X \mid Y\} \,\big\} : \]
the overall mean is the probability-weighted average, over what $Y$ might be, of the
slice-by-slice means.
\end{Obj}
\begin{Prf}{B.2.9}{Proof of Iterated Expectations by Regrouping a Double Sum.}{Uses \xref{B.2.6}, \xref{B.2.8}.}
\item Write out the right side from the definitions. The outer expectation averages the
      function $y\mapsto\E\{X\mid Y=y\}$ against the marginal of $Y$:
      \[ \E\big\{ \E\{X\mid Y\} \big\}
         \;=\; \sum_y \E\{X \mid Y=y\}\; p_Y(y)
         \;=\; \sum_y \Big( \sum_x x\, p(x\mid y) \Big) p_Y(y) . \]
\item Substitute the definition of the conditional pmf,
      $p(x\mid y)=p(x,y)/p_Y(y)$:
      \[ \sum_y \sum_x x\,\frac{p(x,y)}{p_Y(y)}\; p_Y(y)
         \;=\; \sum_y \sum_x x\, p(x,y) , \]
      the marginal $p_Y(y)$ cancelling exactly---this cancellation is the whole proof.
\item Regroup the double sum by $x$ (legal because the sum converges absolutely):
      \[ \sum_x x \sum_y p(x,y) \;=\; \sum_x x\, p_X(x) \;=\; \E\{X\} , \]
      the inner sum being the marginal of $X$ by \xref{B.2.6}.
\end{Prf}

\medskip
Why this theorem matters here: it lets you compute a complicated overall average by
\emph{conditioning}---break the world into slices where things are simple, average within
each slice, then average the slice-answers. The main theorem of \S\,\xref{B.5} is exactly
one application of this maneuver.

% ============================================================
\sub{B.3}{Estimation as a Decision, and Why It Gets Stuck}

The toolkit is assembled; now aim it at the trial of \xref{B.1.1}. This section makes four
ideas precise---model, estimator, loss, risk---and then demonstrates the problem that
motivates everything after: without background knowledge, competing estimators
\emph{cannot be ranked}.

\begin{Obj}{Definition}{B.3.1}{The statistical model, and the trial's}
A \emph{statistical model} is a family of candidate distributions for the data, indexed by
the unknown quantity: for each possible value of $\theta$, a statement of how the data would
then behave. The unknown $\theta$ is called the \emph{parameter}; the set $\Theta$ of its
possible values, the \emph{parameter space}. For the trial: each of $n$ patients is
independently cured with probability $\theta$, and $X$ counts the cures. A particular
sequence with $x$ specified cures and $n-x$ specified failures has probability
$\theta^x(1-\theta)^{n-x}$ (independence multiplies the per-patient probabilities), and
there are $\binom{n}{x}$ ways to choose \emph{which} $x$ patients are the cured ones, so
\[ f(x \mid \theta) \;=\; P(X = x \mid \theta)
   \;=\; \binom{n}{x}\,\theta^{x}(1-\theta)^{\,n-x} ,
   \qquad x = 0, 1, \dots, n , \]
with $\Theta=[0,1]$. This family is the \emph{binomial} model, written
$X\sim\text{Binomial}(n,\theta)$.\note{The notation $f(x\mid\theta)$---``the pmf of the data
\emph{given} the parameter''---is chosen deliberately: from \S\,\xref{B.4} on, $\theta$ will
be treated as a random variable and this really will be a conditional distribution in the
sense of \xref{B.2.6}. The course notes write $f(x;\theta)$ when $\theta$ is being treated
as a fixed unknown; the two notations name the same function.}
\end{Obj}

\begin{Obj}{Definition}{B.3.2}{Estimator}
An \emph{estimator} is a rule---a function $a(\cdot)$, committed to before the data
arrive---that turns any possible data set into a guess: observe $X=x$, report $a(x)$. The
\emph{estimator} is the recipe; the \emph{estimate} is the number the recipe produces on your
data. Candidate recipes for the trial are cheap to write down: report the observed rate
$a_1(x)=x/n$; ignore the data and always report $a_2(x)=\tfrac12$; split the difference,
$a_3(x)=(x+1)/(n+2)$. Each produces \emph{an} answer on every data set, so producing answers
cannot be what makes an estimator good. We need a way to score them.
\end{Obj}

\begin{Obj}{Definition}{B.3.3}{Loss function; squared-error loss}
A \emph{loss function} $L(\theta,a)$ prices the consequence of reporting the guess $a$ when
the truth is $\theta$: bigger loss, worse outcome. It is chosen by \emph{us}, when we pose
the problem, as the definition of what counts as a miss. This handout uses the standard
choice for estimating a number, \emph{squared-error loss}:
\[ L(\theta, a) \;=\; (\theta - a)^2 . \]
Why this and not, say, $|\theta-a|$? Three reasons, in increasing order of importance:
squaring is symmetric (over- and under-estimation are punished alike); it charges large
misses disproportionately (missing by $0.2$ costs four times missing by $0.1$), which
matches how bad decisions usually scale; and---the deep reason, proved in
\S\,\xref{B.5}---it is the loss under which \emph{averages} are the optimal guesses, making
the whole theory computable. Different losses lead to different optimal summaries; squared
error leads to means.\note{For instance, under absolute-error loss $|\theta-a|$ the optimal
guess turns out to be the \emph{median} of one's distribution rather than the mean. The loss
function is not a law of nature; it is the posed question, and it selects the answer.}
\end{Obj}

\begin{Obj}{Definition}{B.3.4}{Risk}
Fix a candidate truth $\theta$ and an estimator $a(\cdot)$. Because the data are random, the
loss you will suffer is random too; its expectation, computed under the model at that
$\theta$, is the \emph{risk}:
\[ R(\theta) \;=\; \E\big\{ L(\theta, a(X)) \big\}
   \;=\; \E\big\{ (\theta - a(X))^2 \big\}
   \qquad\text{($X$ distributed per $f(\cdot\mid\theta)$)} . \]
The risk answers: \emph{if the truth were $\theta$, how badly would this recipe do on
average?} Note carefully that it is a \emph{function} of $\theta$---one number per candidate
truth---because the truth is exactly what we do not know.
\end{Obj}

\begin{ObjL}{Example}{B.3.5}{The risk of the observed rate, computed in full}
Take $a_1(X)=X/n$, the observed rate, under the binomial model. First its mean. Writing
$X=Z_1+\cdots+Z_n$ as the sum of the $n$ per-patient Bernoulli indicators and using
linearity (\xref{B.2.4}):
\[ \E\{X\} \;=\; \E\{Z_1\} + \cdots + \E\{Z_n\} \;=\; \underbrace{\theta+\cdots+\theta}_{n}
   \;=\; n\theta ,
   \qquad\text{so}\qquad
   \E\Big\{\frac{X}{n}\Big\} \;=\; \frac{n\theta}{n} \;=\; \theta . \]
The rate is centered on the truth at \emph{every} $\theta$ (such an estimator is called
\emph{unbiased}). Its risk is therefore its variance:
\[ R_1(\theta) \;=\; \E\Big\{ \Big(\frac{X}{n}-\theta\Big)^{\!2} \Big\}
   \;=\; \E\Big\{ \Big(\frac{X}{n}-\E\Big\{\frac{X}{n}\Big\}\Big)^{\!2} \Big\}
   \;=\; \Var\Big(\frac{X}{n}\Big) , \]
by the definition of variance (\xref{B.2.5}) with the mean $\theta$ substituted. Variances
of independent summands add,\note{Why: for independent $U,V$ with means $\mu_U,\mu_V$, the
cross term in $\E\{((U-\mu_U)+(V-\mu_V))^2\}$ is $2\,\E\{(U-\mu_U)(V-\mu_V)\}
=2\,\E\{U-\mu_U\}\E\{V-\mu_V\}=0$---independence lets the expectation of a product factor
into a product of expectations, and each factor is zero.} and each Bernoulli indicator has
variance $\theta(1-\theta)$ (\xref{B.2.5}), so
\[ \Var(X) \;=\; n\,\theta(1-\theta)
   \qquad\Longrightarrow\qquad
   R_1(\theta) \;=\; \Var\Big(\frac{X}{n}\Big) \;=\; \frac{\Var(X)}{n^2}
   \;=\; \frac{\theta(1-\theta)}{n} . \]
\end{ObjL}

\begin{ObjL}{Example}{B.3.6}{The stuck comparison}
Now score the data-ignoring rule $a_2\equiv\tfrac12$. Its loss is not random at all---it
reports $\tfrac12$ no matter what---so its risk is just
\[ R_2(\theta) \;=\; \big(\theta - \tfrac12\big)^2 . \]
Compare the two at $n=10$. At $\theta=\tfrac12$:
\[ R_1\big(\tfrac12\big) \;=\; \frac{\tfrac12\cdot\tfrac12}{10} \;=\; 0.025 ,
   \qquad
   R_2\big(\tfrac12\big) \;=\; 0 : \]
the stubborn rule is \emph{unbeatable}---it happens to be guessing the truth exactly. But at
$\theta=0.9$:
\[ R_1(0.9) \;=\; \frac{0.9\times0.1}{10} \;=\; 0.009 ,
   \qquad
   R_2(0.9) \;=\; (0.9-0.5)^2 \;=\; 0.16 : \]
now the stubborn rule is nearly eighteen times worse. Neither estimator wins at every
$\theta$; their risk curves \emph{cross}. And this is the general situation, not a quirk of
a silly example: reasonable estimators trade performance in one region of $\Theta$ for
performance in another, and ``minimize the risk'' is not an instruction anyone can follow,
because there is no single risk number to minimize---there is a curve, and curves are not
ordered. The problem is stuck. To get unstuck we must be willing to say \emph{which values
of $\theta$ we care about more}---to put weights on the candidate truths. Doing that with a
probability distribution is the Bayesian move, and it is the next section.
\end{ObjL}

% ============================================================
\sub{B.4}{The Prior, the Posterior, and Bayes' Theorem}

\begin{Obj}{Definition}{B.4.1}{Prior distribution}
A \emph{prior distribution} is a probability distribution $\pi(\theta)$ placed on the
parameter itself, expressing what is believed about $\theta$ \emph{before} the data are
seen: regions of $\Theta$ where $\pi$ is large are values considered plausible, regions
where $\pi$ is small, implausible. Formally this is a change of status: $\theta$ is now
treated as a \emph{random variable} with density (or pmf) $\pi$, so that all of
\S\,\xref{B.2}'s machinery---joints, slices, conditional expectations---applies to it. The
prior is an \emph{input}, chosen by the analyst as a model of background knowledge, exactly
as the binomial model was chosen as a model of the data mechanism. It is not derived from
the data and it is not the thing we will optimize; it is a premise.\note{Two readings
coexist. If $\theta$ was literally produced by a random mechanism (a factory turns out
batches whose defect rates vary), the prior is that mechanism's distribution, an ordinary
frequency statement. More commonly nothing generated $\theta$---the cure rate simply
\emph{is}---and the prior is a quantified statement of evidence and judgment: ``I regard
values near $0.4$ as twice as plausible as values near $0.7$.'' The mathematics downstream
is identical under either reading; the interpretation debate is real and is taken up in
\S\,\xref{B.9}.}
\end{Obj}

\begin{Obj}{Definition}{B.4.2}{The joint distribution of parameter and data}
Once $\theta$ is random, the pair $(\theta, X)$ has a joint distribution, and it is built
from the two models already on the table, prior first, data mechanism second:
\[ \text{joint}(\theta, x) \;=\; \pi(\theta)\; f(x \mid \theta) . \]
Read it generatively: nature draws a truth $\theta$ from $\pi$, then the experiment draws
data $x$ from the model at that truth. (This is the same construction as the table of
\xref{B.2.7}: row weights $0.1/0.9$ were the prior over $D$; the within-row test behavior
was the model; cells were prior $\times$ model.)
\end{Obj}

\begin{Obj}{Definition}{B.4.3}{The marginal distribution of the data}
Summing (or integrating) the joint over $\theta$ gives the \emph{marginal} distribution of
the data,
\[ m(x) \;=\; \int_{\Theta} \pi(\theta)\, f(x\mid\theta)\, d\theta : \]
how probable the data set $x$ is \emph{overall}, averaged over what the truth might be. In
the table it was the column totals. It involves no unknown---$\theta$ has been integrated
away---and its role in what follows is purely that of a normalizing constant.
\end{Obj}

\begin{Obj}{Theorem}{B.4.4}{Bayes' theorem}
The conditional distribution of the parameter given the observed data---the
\emph{posterior distribution}---is
\[ \pi(\theta \mid x) \;=\; \frac{\pi(\theta)\, f(x\mid\theta)}{m(x)}
   \;=\; \frac{\pi(\theta)\, f(x\mid\theta)}
   {\displaystyle\int_{\Theta} \pi(t)\, f(x\mid t)\, dt} . \]
\end{Obj}
\begin{Prf}{B.4.4}{Proof of Bayes' Theorem by Slicing the Joint.}{Uses \xref{B.2.6}, \xref{B.4.2}, \xref{B.4.3}.}
\item By the definition of a conditional distribution (\xref{B.2.6}), applied to the pair
      $(\theta, X)$ with the roles ``condition on the second coordinate'': the conditional
      of $\theta$ given $X=x$ is the joint divided by the marginal of the conditioning
      variable,
      \[ \pi(\theta\mid x) \;=\; \frac{\text{joint}(\theta,x)}{m(x)} . \]
\item Substitute the joint's factorization $\pi(\theta)f(x\mid\theta)$ (\xref{B.4.2}) in
      the numerator and the integral form of $m(x)$ (\xref{B.4.3}) in the denominator.
      (The dummy variable in the denominator is renamed $t$ so it cannot be confused with
      the free $\theta$ in the numerator.)
\end{Prf}

\begin{ObjL}{Remark}{B.4.5}{Reading Bayes' theorem}
Three observations, each load-bearing.
First, the structure is \emph{slice and renormalize}, nothing more: observing $x$ freezes
one slice of the joint; dividing by the slice's total mass $m(x)$ makes it a probability
distribution over $\theta$ again. The theorem is \xref{B.2.6} wearing new labels.
Second, look at how the data enter: only through the factor $f(x\mid\theta)$, read as a
function of $\theta$ with the observed $x$ plugged in. That function has a name---the
\emph{likelihood}---and it scores each candidate truth by how well it explains what was
actually seen. Posterior $\propto$ prior $\times$ likelihood: updated belief is old belief
reweighted by explanatory success.
Third, the denominator does not involve $\theta$ at all. So to identify the posterior it is
enough to know the numerator's \emph{shape} in $\theta$; the constant is whatever makes the
total $1$. In practice one therefore writes
\[ \pi(\theta\mid x) \;\propto\; \pi(\theta)\, f(x\mid\theta) , \]
drops every factor not involving $\theta$, and tries to \emph{recognize} the resulting
shape as a known family. \S\,\xref{B.7} runs exactly this play, and the integral in the
denominator never gets computed.
\end{ObjL}

\begin{ObjL}{Example}{B.4.6}{The table of \xref{B.2.7}, renamed}
You have already done a full Bayesian update. In \xref{B.2.7}: the prior was
$\pi(1)=P(D=1)=0.10$; the model was the test's behavior given the disease state,
$f(1\mid D=1)=0.9$ and $f(1\mid D=0)=0.2$; the observed data was $T=1$; the marginal was
$m(1)=0.10\times0.9+0.90\times0.2=0.27$; and Bayes' theorem gave the posterior
\[ \pi(1\mid 1) \;=\; \frac{\pi(1)\,f(1\mid1)}{m(1)}
   \;=\; \frac{0.10\times0.90}{0.27} \;=\; \frac{0.09}{0.27} \;=\; \frac13 . \]
Belief $0.10$, data, belief $\tfrac13$. Everything from here on is this computation with a
continuous unknown---and then one further question the table never asked: once you have the
posterior, \emph{what single number should you report?}
\end{ObjL}

% ============================================================
\sub{B.5}{The Best Guess Under Squared Error}

The posterior is a whole distribution over $\theta$---a full account of what is now
believed. Often that is the honest thing to report. But estimation demands one number, and
now the loss function chooses it. This section proves the two results on which everything
rests: with no data, the best single-number guess of a random quantity is its \emph{mean};
with data, the best guess is the mean of the \emph{posterior}. The first result is the
engine; the second is one application of iterated expectations to it.

\begin{Obj}{Theorem}{B.5.1}{The best constant guess is the mean}
Let $\theta$ be a random variable with $\E\{\theta^2\}<\infty$. Among all constants $a$,
the expected squared error
\[ \E\big\{ (\theta - a)^2 \big\} \;=\; \Var(\theta) + \big( \E\{\theta\} - a \big)^2 \]
is minimized at exactly one point: $a = \E\{\theta\}$, where it equals $\Var(\theta)$.
\end{Obj}
\begin{Prf}{B.5.1}{Proof of the Best Constant by Adding and Subtracting the Mean.}{Uses \xref{B.2.3}, \xref{B.2.4}, \xref{B.2.5}.}
\item Write $\mu=\E\{\theta\}$ for brevity, and route the error through $\mu$ by adding and
      subtracting it---the step that makes the square split cleanly:
      \[ \theta - a \;=\; (\theta - \mu) + (\mu - a) . \]
\item Square the two-term sum using $(u+v)^2=u^2+2uv+v^2$:
      \[ (\theta-a)^2 \;=\; (\theta-\mu)^2 \;+\; 2(\theta-\mu)(\mu-a) \;+\; (\mu-a)^2 . \]
\item Take expectations term by term (linearity, \xref{B.2.4}). The first term gives
      $\E\{(\theta-\mu)^2\}=\Var(\theta)$ by definition. In the middle term, $\mu-a$ is a
      constant and factors out:
      \[ \E\big\{2(\theta-\mu)(\mu-a)\big\} \;=\; 2(\mu-a)\,\E\{\theta-\mu\}
         \;=\; 2(\mu-a)\big(\E\{\theta\}-\mu\big) \;=\; 2(\mu-a)\cdot 0 \;=\; 0 : \]
      deviations from one's own mean average to zero, so the cross term dies. The last term
      is the constant $(\mu-a)^2$, its own expectation.
\item What remains is the display:
      $\E\{(\theta-a)^2\} = \Var(\theta) + (\mu-a)^2$. The first piece does not involve
      $a$ at all; the second is a square, so it is $\ge 0$ always, $=0$ exactly at
      $a=\mu$, and $>0$ elsewhere. The minimizer is $a=\mu$, uniquely.
\end{Prf}

\begin{ObjL}{Remark}{B.5.2}{What the identity says}
The proof delivered more than a minimizer; it delivered an exact accounting. Guessing a
random quantity with a constant costs you $\Var(\theta)$ \emph{no matter what}---the
irreducible price of the quantity's own randomness---plus a penalty $(\E\{\theta\}-a)^2$
for centering your guess anywhere but the mean. This is why ``balancing point'' was the
right picture for the expectation in \xref{B.2.3}: under squared error, the balancing point
is provably the best single summary. And note the theorem's scope: \emph{any} random
$\theta$, any distribution with a finite second moment. In particular it will apply
verbatim to $\theta$ under its posterior distribution---which is the next theorem.
\end{ObjL}

\begin{Obj}{Definition}{B.5.3}{Bayes risk}
Let $a(\cdot)$ be an estimator. Its \emph{Bayes risk} is its expected squared error when
\emph{both} sources of randomness run: the truth $\theta$ drawn from the prior, the data
$X$ drawn from the model given that truth---that is, $(\theta,X)$ carrying the joint
distribution of \xref{B.4.2}:
\[ r(a) \;=\; \E\big\{ \big( a(X) - \theta \big)^2 \big\} . \]
Compare \xref{B.3.4}: the risk froze $\theta$ and averaged over data only, producing a
curve; the Bayes risk averages over $\theta$ as well---the prior supplies precisely the
weights over candidate truths that \xref{B.3.6} found missing---and produces a \emph{single
number} per estimator. Numbers, unlike curves, can be minimized. The estimator with the
smallest Bayes risk is called the \emph{Bayes estimator}, and the next theorem identifies
it.
\end{Obj}

\begin{Obj}{Theorem}{B.5.4}{The Bayes estimator is the posterior mean}
Assume $\E\{\theta^2\}<\infty$. Among all estimators $a(\cdot)$, the Bayes risk $r(a)$ is
minimized by reporting, for each possible data set $x$, the mean of the posterior
distribution:
\[ a^{*}(x) \;=\; \E\{\theta \mid X = x\}
   \;=\; \int_{\Theta} \theta\; \pi(\theta\mid x)\, d\theta . \]
\end{Obj}
\begin{Prf}{B.5.4}{Proof of the Posterior-Mean Rule by Conditioning on the Data.}{Uses \xref{B.2.9}, \xref{B.4.4}, \xref{B.5.1}.}
\item Iterate the Bayes risk through the data, using iterated expectations (\xref{B.2.9})
      with the conditioning variable $X$:
      \[ r(a) \;=\; \E\big\{ (a(X)-\theta)^2 \big\}
         \;=\; \E\Big[\, \E\big\{ (a(X)-\theta)^2 \bigm| X \big\} \,\Big] : \]
      the overall average equals the average, over data sets, of the within-data-set
      averages.
\item Freeze one slice $X=x$ and look at the inner expectation. Inside the slice, $a(x)$ is
      a fixed \emph{constant}---the recipe's output on that data set---while $\theta$
      remains random, distributed according to the conditional distribution of $\theta$
      given $X=x$. And that conditional distribution is, by \xref{B.4.4}, exactly the
      posterior $\pi(\cdot\mid x)$. So the inner expectation is the expected squared error
      of a \emph{constant} guess at a random quantity:
      \[ \E\big\{ (a(x) - \theta)^2 \bigm| X=x \big\}
         \;=\; \E\big\{ (\theta - c)^2 \big\}\Big|_{c=a(x)}
         \quad\text{with $\theta$ carrying the posterior.} \]
\item Apply \xref{B.5.1} \emph{to the posterior distribution}---the theorem holds for any
      distribution, so in particular this one. Within the slice, the inner expectation is
      minimized, uniquely, by the constant equal to that distribution's mean:
      \[ c \;=\; \E\{\theta \mid X = x\} , \]
      and no other choice of $a(x)$ does as well in this slice.
\item The rule $a^{*}(x)=\E\{\theta\mid X=x\}$ therefore minimizes the inner expectation
      \emph{simultaneously at every} $x$. Averaging over slices (step (1)) preserves the
      domination: for any competitor $a$, each of $a$'s slice-averages is at least
      $a^{*}$'s, so the overall average is too. Hence $r(a^{*})\le r(a)$ for every
      estimator $a$, which is the claim.
\end{Prf}

\begin{ObjL}{Remark}{B.5.5}{The recipe, in words}
The theorem compresses the whole method into two verbs: \textbf{update, then average}.
\begin{enumerate}[leftmargin=2.2em,itemsep=2pt,topsep=2pt]
  \item Model the data mechanism ($f(x\mid\theta)$) and state your prior ($\pi$).
  \item Observe $x$; form the posterior $\pi(\theta\mid x)\propto\pi(\theta)f(x\mid\theta)$
        (\xref{B.4.4}), the updated belief.
  \item Report the posterior's mean. That single number minimizes expected squared error
        among \emph{all} possible estimators, given the premises.
\end{enumerate}
Notice what is chosen and what is given: the prior and the loss are premises, posed with the
problem; the theorem then leaves \emph{nothing} to choose---the optimal estimator is
determined. The remaining work in any particular problem is computational: actually finding
the posterior and its mean. The next three sections do that work twice, in the two settings
where it is fully explicit.
\end{ObjL}

% ============================================================
\sub{B.6}{A Toolkit Distribution: the Beta Family}

To run the recipe on the trial we need a family of prior distributions for a
\emph{fraction}---a random quantity living in $[0,1]$---rich enough to express real beliefs
and tame enough to compute with. The standard choice is the Beta family. Its normalizing
constant involves the Gamma function, so that comes first; both objects repay the ten
minutes they take to set up.

\begin{Obj}{Definition}{B.6.1}{The Gamma function}
For $s>0$,
\[ \Gamma(s) \;=\; \int_0^{\infty} t^{\,s-1}\, e^{-t}\, dt . \]
Two facts are all we need. First, the recursion $\Gamma(s+1)=s\,\Gamma(s)$: integrate by
parts with $u=t^{s}$ (so $du=s\,t^{s-1}dt$) and $dv=e^{-t}dt$ (so $v=-e^{-t}$):
\[ \Gamma(s+1) \;=\; \int_0^{\infty} t^{s} e^{-t}\, dt
   \;=\; \Big[ -t^{s} e^{-t} \Big]_0^{\infty} + s\int_0^{\infty} t^{\,s-1} e^{-t}\, dt
   \;=\; 0 + s\,\Gamma(s) , \]
the boundary term vanishing at both ends ($t^s e^{-t}\to0$ as $t\to\infty$ because the
exponential beats any power; at $t=0$, $t^s=0$ for $s>0$). Second, the anchor
$\Gamma(1)=\int_0^\infty e^{-t}dt=1$. Together they give, for whole numbers,
$\Gamma(n) = (n-1)\,\Gamma(n-1) = \cdots = (n-1)!$\,: the Gamma function is the factorial,
shifted by one and extended to non-integers.
\end{Obj}

\begin{Obj}{Definition}{B.6.2}{The Beta distribution}
For parameters $\alpha>0$, $\beta>0$, the \emph{Beta$(\alpha,\beta)$ distribution} is the
continuous distribution on $[0,1]$ with density
\[ \pi(\theta) \;=\; \frac{\theta^{\,\alpha-1}(1-\theta)^{\,\beta-1}}{\mathrm{B}(\alpha,\beta)} ,
   \qquad
   \mathrm{B}(\alpha,\beta) \;=\; \frac{\Gamma(\alpha)\,\Gamma(\beta)}{\Gamma(\alpha+\beta)} , \]
the constant $\mathrm{B}(\alpha,\beta)$ being exactly what makes the density integrate to
$1$. The exponents steer the shape: mass shifts toward $1$ as $\alpha$ grows and toward $0$
as $\beta$ grows, so $(\alpha,\beta)$ behave like counts of imagined prior successes and
failures---a reading \xref{B.7.3} will make exact. One member deserves special mention:
$\alpha=\beta=1$ gives $\pi(\theta)=\theta^{0}(1-\theta)^{0}/\mathrm{B}(1,1)$, and since
$\mathrm{B}(1,1)=\Gamma(1)\Gamma(1)/\Gamma(2)=1\cdot1/1=1$, this is $\pi(\theta)\equiv1$:
the \emph{uniform} distribution, the flat prior that treats every fraction as equally
plausible.
\end{Obj}

\begin{Obj}{Theorem}{B.6.3}{The mean of a Beta distribution}
If $\theta\sim$ Beta$(\alpha,\beta)$, then
\[ \E\{\theta\} \;=\; \frac{\alpha}{\alpha+\beta} . \]
\end{Obj}
\begin{Prf}{B.6.3}{Proof of the Beta Mean by the Kernel Trick.}{Uses \xref{B.2.3}, \xref{B.6.1}, \xref{B.6.2}.}
\item From the definition of expectation, multiplying the density by $\theta$ merges the
      extra $\theta$ into the exponent:
      \[ \E\{\theta\} \;=\; \int_0^1 \theta\cdot
         \frac{\theta^{\,\alpha-1}(1-\theta)^{\,\beta-1}}{\mathrm{B}(\alpha,\beta)}\,d\theta
         \;=\; \frac{1}{\mathrm{B}(\alpha,\beta)}
         \int_0^1 \theta^{\,(\alpha+1)-1}(1-\theta)^{\,\beta-1}\,d\theta . \]
\item Recognize the remaining integrand as the \emph{un}normalized Beta$(\alpha+1,\beta)$
      density. Its integral must therefore be that family's normalizing constant:
      \[ \int_0^1 \theta^{\,(\alpha+1)-1}(1-\theta)^{\,\beta-1}\,d\theta
         \;=\; \mathrm{B}(\alpha+1,\beta) . \]
      (This maneuver---evaluate an integral by recognizing a known density's shape---is the
      \emph{kernel trick}, and it is the workhorse of every Bayesian computation to come.)
\item So $\E\{\theta\}=\mathrm{B}(\alpha+1,\beta)/\mathrm{B}(\alpha,\beta)$. Expand both
      via \xref{B.6.2} and cancel:
      \[ \E\{\theta\}
         \;=\; \frac{\Gamma(\alpha+1)\,\Gamma(\beta)}{\Gamma(\alpha+\beta+1)}\cdot
               \frac{\Gamma(\alpha+\beta)}{\Gamma(\alpha)\,\Gamma(\beta)}
         \;=\; \frac{\Gamma(\alpha+1)}{\Gamma(\alpha)}\cdot
               \frac{\Gamma(\alpha+\beta)}{\Gamma(\alpha+\beta+1)} . \]
\item Apply the recursion $\Gamma(s+1)=s\,\Gamma(s)$ (\xref{B.6.1}) to each ratio:
      $\Gamma(\alpha+1)/\Gamma(\alpha)=\alpha$ and
      $\Gamma(\alpha+\beta)/\Gamma(\alpha+\beta+1)=1/(\alpha+\beta)$. Hence
      $\E\{\theta\}=\alpha/(\alpha+\beta)$.
\end{Prf}

% ============================================================
\sub{B.7}{The Cure Rate, Start to Finish}

Everything is on the table. Model: $X\sim$ Binomial$(n,\theta)$ (\xref{B.3.1}). Prior: to
begin, the flat Beta$(1,1)$---``before the trial, no fraction is privileged over any
other.'' Loss: squared error. Data: $x=7$ cures among $n=10$ patients. The recipe of
\xref{B.5.5} says: form the posterior, report its mean.

\begin{ObjL}{Example}{B.7.1}{The posterior, by kernel-matching}
Bayes' theorem in its proportional form (\xref{B.4.5}), for a general $x$ and $n$ under the
flat prior:
\[ \pi(\theta \mid x) \;\propto\; \pi(\theta)\, f(x\mid\theta)
   \;=\; 1 \cdot \binom{n}{x}\,\theta^{x}(1-\theta)^{\,n-x}
   \;\propto\; \theta^{x}(1-\theta)^{\,n-x} , \]
the binomial coefficient dropping because it does not involve $\theta$---it is data-only,
hence absorbed into the normalizing constant. Now \emph{match the kernel}: rewrite the
exponents in Beta form,
\[ \theta^{x}(1-\theta)^{\,n-x} \;=\; \theta^{\,(x+1)-1}\,(1-\theta)^{\,(n-x+1)-1} , \]
which is precisely the shape of a Beta$(x+1,\; n-x+1)$ density (\xref{B.6.2}). A density's
shape determines it (the constant is forced by total mass $1$), so
\[ \theta \mid X = x \;\sim\; \text{Beta}\big(x+1,\; n-x+1\big) , \]
and the normalizing integral $m(x)$ of \xref{B.4.3} was never touched. By \xref{B.6.3}, the
Bayes estimate---the posterior mean---is
\[ a^{*}(x) \;=\; \E\{\theta\mid X=x\} \;=\; \frac{x+1}{(x+1)+(n-x+1)}
   \;=\; \frac{x+1}{n+2} . \]
\end{ObjL}

\begin{ObjL}{Example}{B.7.2}{The trial's numbers}
With $n=10$ and $x=7$: the posterior is Beta$(8,4)$, and the Bayes estimate is
\[ a^{*}(7) \;=\; \frac{7+1}{10+2} \;=\; \frac{8}{12} \;=\; \frac{2}{3} \;\approx\; 0.667 , \]
against the raw rate's $x/n=0.700$. The posterior also prices its own uncertainty: the
Beta$(\alpha,\beta)$ variance is
$\alpha\beta/\big((\alpha+\beta)^2(\alpha+\beta+1)\big)$,\note{Derived by the same kernel
trick as the mean: $\E\{\theta^2\}=\mathrm{B}(\alpha+2,\beta)/\mathrm{B}(\alpha,\beta)
=\frac{(\alpha+1)\alpha}{(\alpha+\beta+1)(\alpha+\beta)}$ (two applications of the
recursion), then $\Var=\E\{\theta^2\}-\E\{\theta\}^2$ with common denominator
$(\alpha+\beta)^2(\alpha+\beta+1)$.} here
\[ \Var(\theta\mid X=7) \;=\; \frac{8\cdot4}{12^2\cdot13} \;=\; \frac{32}{1872}
   \;=\; \frac{2}{117} \;\approx\; 0.0171 ,
   \qquad \text{sd} \;\approx\; 0.131 . \]
So the full Bayesian report is: \emph{the cure rate is around $0.67$, give or take about
$0.13$}---an estimate and an honest width, both read off one distribution.
\end{ObjL}

\begin{ObjL}{Remark}{B.7.3}{The anatomy of $(x+1)/(n+2)$: a weighted average}
Rewrite the Bayes estimate to expose its structure:
\[ \frac{x+1}{n+2}
   \;=\; \frac{2\cdot\tfrac12 \;+\; n\cdot\dfrac{x}{n}}{2+n} . \]
(Check: the numerator is $1+x$ and the denominator $2+n$, as before.) Read it as a weighted
average of two guesses: the \emph{prior mean} $\tfrac12$, carrying weight $2$, and the
\emph{data's rate} $x/n$, carrying weight $n$. The flat prior behaves exactly like
\emph{two phantom patients, one cured and one not}, mixed into the sample. The same algebra
runs for any Beta$(\alpha,\beta)$ prior: the posterior is
Beta$(\alpha+x,\,\beta+n-x)$---observed successes reinforce $\alpha$, failures reinforce
$\beta$---and its mean is
\[ \frac{\alpha+x}{\alpha+\beta+n}
   \;=\; \frac{(\alpha+\beta)\cdot\dfrac{\alpha}{\alpha+\beta} \;+\; n\cdot\dfrac{x}{n}}
   {(\alpha+\beta)+n} : \]
prior mean weighted by $\alpha+\beta$ (\emph{the prior is worth $\alpha+\beta$
observations}), data rate weighted by $n$. This is the general shape of Bayesian
estimates: a compromise between what you believed and what you saw, at exchange rates set
by how much each side knows.
\end{ObjL}

\begin{Fig}{B.7.4}{The update, drawn}{The flat prior (dashed) and the posterior
Beta$(8,4)$ after observing $7$ cures in $10$ patients. The posterior's \emph{peak} sits at
exactly the raw rate $x/n=0.7$ (the data's best single explanation), but its \emph{mean}---%
the Bayes estimate under squared error, marked in green---sits at $2/3\approx0.667$, pulled
toward the middle by the distribution's left tail. The whole curve, not just one number, is
the updated state of knowledge; the loss function then selects which one number to report
(\xref{B.5.1}).}
\begin{tikzpicture}[x=7cm,y=0.85cm, every node/.style={font=\small}]
  % axis
  \draw[black!60] (-0.02,0) -- (1.04,0);
  \foreach \t/\l in {0/0, 0.25/0.25, 0.5/0.5, 0.75/0.75, 1/1}{
    \draw[black!60] (\t,0) -- (\t,-0.12);
    \node[black!70, font=\footnotesize, below] at (\t,-0.10) {$\l$};
  }
  \node[black!70, font=\footnotesize, below] at (1.065,-0.08) {$\theta$};
  % flat prior
  \draw[black!55, densely dashed, thick] (0,1) -- (1,1);
  \node[black!55, font=\footnotesize, above] at (0.13,1.02) {prior Beta$(1,1)$};
  % posterior Beta(8,4): f = 1320 t^7 (1-t)^3
  \draw[figTerm, very thick, domain=0.18:0.995, smooth, samples=120]
    plot (\x, {1320*(\x)^7*(1-\x)^3});
  \node[figTerm, font=\footnotesize, above left] at (0.56,2.15) {posterior Beta$(8,4)$};
  % posterior mean 2/3
  \draw[figLim, densely dashed, thick] (0.6667,0) -- (0.6667,{1320*0.6667^7*0.3333^3});
  \node[figLim!80!black, font=\footnotesize, below] at (0.645,-0.55)
    {$\E\{\theta\mid x\}=\tfrac23$};
  % mode = raw rate 0.7
  \fill[figTerm] (0.7,{1320*0.7^7*0.3^3}) circle (1.4pt);
  \node[figTerm, font=\footnotesize, above] at (0.79,{1320*0.7^7*0.3^3+0.06})
    {mode $=x/n=0.7$};
\end{tikzpicture}
\end{Fig}

\begin{ObjL}{Remark}{B.7.5}{Two sanity checks}
\emph{No data.} Set $n=x=0$: the estimate is $(0+1)/(0+2)=\tfrac12$, the prior mean---with
nothing observed, the recipe reports what you believed. \emph{All cures.} Set $x=n$: the
estimate is $(n+1)/(n+2)$, which is large but strictly less than $1$: ten out of ten cures
should not convince anyone that failure is \emph{impossible}. (The raw rate $x/n=1$ would
assert exactly that.) This estimator for the all-successes case is old and famous---it is
\emph{Laplace's rule of succession}, his answer to ``what probability that the sun rises
tomorrow, given that it always has?''---and it falls out of the machinery here as a routine
special case.
\end{ObjL}

% ============================================================
\sub{B.8}{A Second Instrument: the Normal--Normal Pair}

The Beta--binomial pair handled a fraction. The other workhorse setting is a
\emph{measurement}: a quantity on the whole line, observed through noise. One new
distribution is needed.

\begin{Obj}{Definition}{B.8.1}{The normal distribution}
The \emph{normal} (or Gaussian) distribution $N(\mu,\sigma^2)$ is the continuous
distribution with density
\[ f(y) \;=\; \frac{1}{\sigma\sqrt{2\pi}}\,
   \exp\!\left( -\frac{(y-\mu)^2}{2\sigma^2} \right) , \qquad y\in\mathbb{R} : \]
the bell curve, centered at its mean $\mu$, with spread set by its standard deviation
$\sigma$ (about $68\%$ of the probability lies within one $\sigma$ of $\mu$, about $95\%$
within two). Its ubiquity has a theorem behind it (sums of many small independent effects
are approximately normal---the central limit theorem, not needed here) and a practical
reason: as a function of the variable in the exponent, the density is
$e^{-\text{quadratic}}$, and quadratics are closed under addition. That closure is exactly
what the next example exploits.
\end{Obj}

\begin{ObjL}{Example}{B.8.2}{Normal prior, normal measurement: the posterior in full}
\emph{Setup.} A quantity $\theta$ (an aptitude score, a physical constant, a true weight)
has prior $\theta\sim N(m,\tau^2)$: best prior guess $m$, prior uncertainty $\tau$. An
instrument measures it once with known noise level $\sigma$: given the truth,
$X\mid\theta\sim N(\theta,\sigma^2)$. Observe $X=x$. What is the posterior?
\emph{Multiply prior by likelihood} (\xref{B.4.5}), keeping only factors involving
$\theta$---the two normalizing constants $1/(\tau\sqrt{2\pi})$ and $1/(\sigma\sqrt{2\pi})$
drop:
\[ \pi(\theta\mid x) \;\propto\;
   \exp\!\left( -\frac{(\theta-m)^2}{2\tau^2} \right)\cdot
   \exp\!\left( -\frac{(x-\theta)^2}{2\sigma^2} \right)
   \;=\; \exp\!\left( -\frac12\left[ \frac{(\theta-m)^2}{\tau^2}
   + \frac{(\theta-x)^2}{\sigma^2} \right] \right) . \]
\emph{Expand the bracket and collect powers of $\theta$.} Expanding each square,
\[ \frac{\theta^2 - 2m\theta + m^2}{\tau^2} + \frac{\theta^2 - 2x\theta + x^2}{\sigma^2}
   \;=\; \theta^2\Big( \frac{1}{\tau^2}+\frac{1}{\sigma^2} \Big)
   \;-\; 2\theta\Big( \frac{m}{\tau^2}+\frac{x}{\sigma^2} \Big)
   \;+\; \Big( \frac{m^2}{\tau^2}+\frac{x^2}{\sigma^2} \Big) , \]
and the last group is free of $\theta$---bound for the normalizing constant.
\emph{Complete the square.} Abbreviate
\[ \frac1v \;=\; \frac{1}{\tau^2}+\frac{1}{\sigma^2} , \qquad
   \frac{m^{*}}{v} \;=\; \frac{m}{\tau^2}+\frac{x}{\sigma^2}
   \quad\text{(defining $v$ and $m^*$)} , \]
so the $\theta$-dependent part is
\[ \frac{\theta^2 - 2 m^{*}\theta}{v}
   \;=\; \frac{(\theta - m^{*})^2 - (m^{*})^2}{v} , \]
by adding and subtracting $(m^*)^2$ inside---and $(m^*)^2/v$, constant in $\theta$, also
departs for the normalizer. What remains is
\[ \pi(\theta\mid x) \;\propto\; \exp\!\left( -\frac{(\theta-m^{*})^2}{2v} \right) : \]
the kernel of a normal density (\xref{B.8.1}). Matching shapes,
\[ \theta \mid X=x \;\sim\; N\big( m^{*},\, v \big) , \qquad
   m^{*} \;=\; v\Big( \frac{m}{\tau^2}+\frac{x}{\sigma^2} \Big)
   \;=\; \frac{\sigma^2\, m + \tau^2\, x}{\sigma^2+\tau^2} , \]
the last simplification by multiplying numerator and denominator by $\sigma^2\tau^2$:
\[ v\Big(\frac{m}{\tau^2}+\frac{x}{\sigma^2}\Big)
   = \frac{\frac{m}{\tau^2}+\frac{x}{\sigma^2}}{\frac{1}{\tau^2}+\frac{1}{\sigma^2}}
   = \frac{\sigma^2 m + \tau^2 x}{\sigma^2 + \tau^2} . \]
By \xref{B.5.4} the Bayes estimate is $m^{*}$, and the posterior variance $v$ prices the
remaining uncertainty.
\end{ObjL}

\begin{ObjL}{Example}{B.8.3}{Numbers: one loud measurement}
Prior: $\theta\sim N(100, 15^2)$ (say, an aptitude score known to average $100$ with
spread $15$ in the population the subject comes from). Instrument: $\sigma=10$. Observation:
$x=130$---a strikingly high reading. Posterior, by \xref{B.8.2}, in exact fractions:
\[ \frac1v \;=\; \frac{1}{225}+\frac{1}{100} \;=\; \frac{4}{900}+\frac{9}{900}
   \;=\; \frac{13}{900}
   \qquad\Longrightarrow\qquad v \;=\; \frac{900}{13} \;\approx\; 69.2 ,
   \quad \sqrt{v} \;=\; \frac{30}{\sqrt{13}} \;\approx\; 8.3 ; \]
\[ m^{*} \;=\; v\Big( \frac{100}{225} + \frac{130}{100} \Big)
   \;=\; \frac{900}{13}\Big( \frac{4}{9} + \frac{13}{10} \Big)
   \;=\; \frac{900}{13}\cdot\frac{40 + 117}{90}
   \;=\; \frac{10\cdot 157}{13}
   \;=\; \frac{1570}{13} \;\approx\; 120.8 . \]
The Bayes estimate is about $120.8$: well above the prior's $100$ (the measurement was loud
and the instrument decent), but pulled down meaningfully from the raw reading $130$ (a
single reading, from a population where $130$ is rare, is partly signal and partly luck).
The posterior spread $8.3$ is smaller than \emph{both} the prior's $15$ and the
instrument's $10$: two imperfect sources of information, combined, know more than either
alone.
\end{ObjL}

\begin{Fig}{B.8.4}{Prior, likelihood, posterior}{The update of \xref{B.8.3}: prior
$N(100,15^2)$ (dashed), the likelihood of the observation $x=130$ (blue, a normal shape in
$\theta$ centered at the reading), and the posterior $N(120.8,\,8.3^2)$ (green). The
posterior sits between its two parents, closer to the likelihood because the instrument
($\sigma=10$) is more precise than the prior ($\tau=15$)---and it is narrower than either,
because precisions add (\xref{B.8.5}). Curves drawn to a common vertical scale.}
\begin{tikzpicture}[x=0.075cm,y=1cm, every node/.style={font=\small}]
  % axis
  \draw[black!60] (60,0) -- (162,0);
  \foreach \t in {70, 100, 130}{
    \draw[black!60] (\t,0) -- (\t,-0.10);
    \node[black!70, font=\footnotesize, below] at (\t,-0.08) {$\t$};
  }
  \node[black!70, font=\footnotesize, below] at (165,-0.06) {$\theta$};
  % prior N(100,15^2): 2.394*exp(-(x-100)^2/450)
  \draw[black!55, densely dashed, thick, domain=62:140, smooth, samples=90]
    plot (\x, {2.394*exp(-(\x-100)^2/450)});
  \node[black!55, font=\footnotesize] at (80,2.62) {prior};
  % likelihood N(130,10^2): 3.590*exp(-(x-130)^2/200)
  \draw[figTerm, thick, domain=100:160, smooth, samples=90]
    plot (\x, {3.590*exp(-(\x-130)^2/200)});
  \node[figTerm, font=\footnotesize] at (146,3.45) {likelihood};
  % posterior N(120.77, 69.23): 4.315*exp(-(x-120.77)^2/138.46)
  \draw[figLim, very thick, domain=95:147, smooth, samples=90]
    plot (\x, {4.315*exp(-(\x-120.77)^2/138.46)});
  \node[figLim!85!black, font=\footnotesize] at (110.5,4.35) {posterior};
  % posterior mean tick
  \draw[figLim, densely dashed] (120.77,0) -- (120.77,4.315);
  \node[figLim!80!black, font=\footnotesize, below] at (120.77,-0.08) {$120.8$};
\end{tikzpicture}
\end{Fig}

\begin{ObjL}{Remark}{B.8.5}{The precision reading, and shrinkage}
State \xref{B.8.2}'s formulas in terms of \emph{precision} $=$ $1/\text{variance}$ (how
much a source of information knows):
\[ \underbrace{\frac1v}_{\text{posterior precision}}
   \;=\; \underbrace{\frac{1}{\tau^2}}_{\text{prior precision}}
   + \underbrace{\frac{1}{\sigma^2}}_{\text{data precision}} ,
   \qquad
   m^{*} \;=\; \frac{\frac{1}{\tau^2}}{\frac{1}{\tau^2}+\frac{1}{\sigma^2}}\; m
   \;+\; \frac{\frac{1}{\sigma^2}}{\frac{1}{\tau^2}+\frac{1}{\sigma^2}}\; x : \]
\emph{precisions add}, and the posterior mean weighs each guess by its share of the total
precision---the continuous twin of \xref{B.7.3}'s observation-counting. In the numbers of
\xref{B.8.3} the weights are $\frac{4}{13}$ on the prior and $\frac{9}{13}$ on the data:
$m^*=\frac{4}{13}\cdot100+\frac{9}{13}\cdot130=\frac{400+1170}{13}=\frac{1570}{13}$, as
before. Equivalently,
\[ m^{*} \;=\; m + \frac{\tau^2}{\sigma^2+\tau^2}\,(x - m)
   \;=\; 100 + \frac{225}{325}\,(130-100) \;=\; 100 + \frac{9}{13}\cdot30 \;\approx\;
   120.8 : \]
the estimate starts at the prior mean and covers only the fraction
$\tau^2/(\sigma^2+\tau^2)$ of the distance to the data. Bayesian estimates \emph{shrink}
noisy observations toward what was already believed---gently when the data are precise,
firmly when they are not. This one word, \emph{shrinkage}, is the modern doorway from this
handout into a large part of contemporary statistics and machine learning.
\end{ObjL}

% ============================================================
\sub{B.9}{What to Make of the Prior}

\begin{ObjL}{Remark}{B.9.1}{The estimate depends on the prior---as it must}
Rerun the trial of \xref{B.7.2} under a different premise. Suppose historical experience
with similar treatments made low cure rates plausible---say a Beta$(2,6)$ prior, mean
$2/8=0.25$, worth $\alpha+\beta=8$ phantom observations. By \xref{B.7.3} the posterior
after $7$ cures in $10$ is Beta$(2+7,\,6+3)=$ Beta$(9,9)$, with mean
\[ \frac{2+7}{2+6+10} \;=\; \frac{9}{18} \;=\; 0.500 , \]
against $0.667$ under the flat prior and $0.700$ raw. Ten patients cannot overrule eight
phantoms. This is not a defect; it is the premise doing exactly what it claimed: with
little data, background knowledge \emph{should} matter, and the machinery is transparent
about how much. But it lays an obligation on the analyst---the prior is part of the
argument, must be stated, and must be defensible on grounds \emph{outside} the data being
analyzed. Estimating the prior's location from the same data it will then shrink is
circular, and the circle does not close.\note{The course deck's aphorism, quoted in the
notes at \xrefL{2.10.2}: ``you cannot use the data to know what to shrink towards; but you
can use the data to know how much to shrink towards it''---the \emph{amount} of shrinkage,
unlike the \emph{target}, is learnable when many parallel problems share one prior.}
\end{ObjL}

\begin{ObjL}{Remark}{B.9.2}{The data win in the end}
The dependence on the prior is a small-sample phenomenon. Compare the Bayes estimate with
the observed rate at stage $n$ by putting them over a common denominator:
\[ \frac{\alpha+x}{\alpha+\beta+n} - \frac{x}{n}
   \;=\; \frac{n(\alpha+x) - x(\alpha+\beta+n)}{n\,(\alpha+\beta+n)}
   \;=\; \frac{\alpha n - x(\alpha+\beta)}{n\,(\alpha+\beta+n)}
   \;=\; \frac{\alpha - (\alpha+\beta)\,\frac{x}{n}}{\alpha+\beta+n} , \]
the middle step cancelling the $xn$ terms. Since the rate $x/n$ lies in $[0,1]$, the
numerator can be no larger in size than $\max(\alpha,\beta)$---a fixed number---while the
denominator grows without bound. So the gap between the Bayes estimate and the observed
rate shrinks to zero as data accumulate, \emph{whatever} the data do: every fixed Beta
prior's estimate is eventually indistinguishable from the raw rate---the phantoms are
outvoted. Priors differing in judgment but agreed on possibilities are eventually
reconciled by evidence; where the prior genuinely matters is exactly where matters are
genuinely uncertain: small samples.
\end{ObjL}

\begin{ObjL}{Remark}{B.9.3}{Two schools, one course}
A statistician who declines to put a distribution on $\theta$---who insists the cure rate
is a fixed fact, not a random variable---works instead with the risk \emph{curves} of
\xref{B.3.4} and principles for comparing them (ruling out estimators beaten everywhere;
guarding the worst case). That is the \emph{frequentist} school, and it is where the course
notes spend \S\,\xrefL{1.5}. The Bayesian school, this handout's subject, buys a complete
and unique answer at the price of a prior. The two are not enemies so much as different
contracts with the reader about what may be assumed---and they meet constantly: the
course's Lecture~1 appendix shows a Bayesian-derived rule, $(x+1)/(n+2)$, holding its own
in a purely frequentist risk comparison (\xrefL{1.A.4}, \xrefL{1.A.5}). A working
statistician reads both languages.
\end{ObjL}

% ============================================================
\sub{B.10}{The Vocabulary, and Where to Go Next}

\begin{ObjL}{Concept}{B.10.1}{The names of everything you just did}
For searching, reading, and talking to statisticians---the standard names, each pointing at
its object in this handout:
\begin{itemize}[leftmargin=1.8em,itemsep=2pt,topsep=2pt]
  \item \emph{Prior distribution} (\xref{B.4.1}); \emph{likelihood} (\xref{B.4.5});
        \emph{posterior distribution} (\xref{B.4.4}); \emph{Bayes' theorem} (\xref{B.4.4});
        \emph{Bayesian updating} (the whole \S\,\xref{B.4} pipeline).
  \item \emph{Bayes risk} (\xref{B.5.3}); \emph{Bayes estimator} or \emph{Bayes rule}
        (\xref{B.5.4})---beware that ``Bayes' rule'' also names the updating theorem;
        context decides. The posterior-mean result is often quoted as ``the Bayes estimator
        under squared-error (quadratic) loss is the posterior mean.''
  \item \emph{Conjugate prior}: a family, like Beta for the binomial (\S\,\xref{B.7}) or
        normal for the normal (\S\,\xref{B.8}), whose posteriors stay in the family---the
        kernel trick succeeds by design. \emph{Shrinkage} (\xref{B.8.5});
        \emph{Laplace's rule of succession} (\xref{B.7.5}).
  \item In engineering literature the \S\,\xref{B.5} mathematics runs under \emph{minimum
        mean squared error (MMSE) estimation}: the conditional expectation as the MMSE
        predictor.
  \item One report this handout did not build: an \emph{interval} from the posterior (e.g.,
        the central $95\%$ of Beta$(8,4)$), called a \emph{credible interval}---the
        Bayesian counterpart of the confidence sets in the course notes,
        \S\,\xrefL{2.7}.
\end{itemize}
\end{ObjL}

\begin{ObjL}{Concept}{B.10.2}{Where to go next}
In the course notes: \S\,\xrefL{2.8} proves the two prediction lemmas in their general form,
\S\,\xrefL{2.9} is this handout's material at course pace (with the Poisson--gamma pair as
a third conjugate example), and \xrefL{1.A.5} runs the trial's computation inside the
decision-theory frame of Lecture~1. In the texts: Hoel--Port--Stone treat Bayesian methods
in \S 2.8 (their Theorem~2 is \xref{B.5.4}); DeGroot \& Schervish devote Chapter~7
(\S\S 7.2--7.4) to priors, posteriors, and Bayes estimators at a gentle pace with many
worked examples; Blitzstein \& Hwang's Chapters 7--9 are the friendliest place to firm up
the conditional-probability foundations of \S\,\xref{B.2}. And carry out one exercise by
hand before reading further: redo \S\,\xref{B.7} with your own prior---pick $\alpha$ and
$\beta$ to express something you actually believe about some fraction in your life, invent
data, and push it through \xref{B.7.3}. The method fits on an index card; the point of the
theory was to earn it.
\end{ObjL}

\ifSubfilesClassLoaded{\clearpage\objindex}{}

\end{document}
